\begin{tikzpicture}
\path[use as bounding box] (-.1,-.7) rectangle (22.0,15.1);
\node[heading] at (11,14.6) {同一个恢复入口，两种保留范围};
\node[note,anchor=west] at (.75,13.8) {两条独立执行都从 $R=115,\ G=108,\ G'=D=105$ 开始；高度是压缩示意，不按秒数换算。};
\node[subheading,text=blue,anchor=west] at (.75,12.95) {A. Espresso 长停机：保住原来的全部 $G$，丢弃尚未进入 $G$ 的区域后缀};
\node[box,kept,fill=settled,text width=3.3cm] (a0) at (2.7,10.1) {\textbf{共同前缀至 108}\\已有 106--108 交付成功};
\node[box,dead,text width=4.0cm] (aold) at (8.8,11.7) {\textbf{旧区域 109--115}\\A 的尾部与 B 的后继};
\node[note,anchor=west,text width=6.5cm] at (12.1,11.7) {窗口回退后，旧 session 被 Kill。\\这些块没有进入旧 $G$。};
\node[box,fill=native,draw=violet,text width=3.4cm] (a1) at (8.8,8.85) {\textbf{$D$：109--120}\\窗口到期，deposit-only};
\node[box,fill=forceB,draw=rose,text width=3.4cm] (a2) at (14.3,8.85) {\textbf{$D$：121--125}\\RESUME(125) 立即补齐};
\node[box,kept,fill=settled,text width=2.6cm] (a3) at (19.6,8.85) {\textbf{新 session}\\从 $D[125]$ 出 126};
\draw[ref] (a0.north east) -- (5.5,11.7) -- (aold.west);
\draw[flow,draw=violet] (a0.south east) -- (5.5,8.85) -- (a1.west);
\draw[recover] (a1) -- (a2); \draw[recover] (a2) -- (a3);
\node[note,anchor=west,text=green] at (.75,7.45) {RESUME 后：$G$ 保留至 108，再采纳 $D$ 的 109--125；原 $G$ 无需撤销。};

\draw[line!35] (.75,6.9) -- (21.3,6.9);
\node[subheading,text=orange,anchor=west] at (.75,6.25) {B. 批次层超过窗口：$D$ 在 106 分叉，旧 $G$ 的未结算后缀被撤销};
\node[box,kept,fill=settled,text width=3.3cm] (b0) at (2.7,3.4) {\textbf{共同前缀至 105}\\Inbox 停在 cursor 105};
\node[box,dead,text width=4.0cm] (bold) at (8.8,5.0) {\textbf{旧 $G$：106--200}\\有全局确认，未按时交付};
\node[note,anchor=west,text=rose,text width=6.5cm] at (12.1,5.0) {Kill 停止旧生产；\\RESUME 才撤销这段 $G$。};
\node[box,fill=native,draw=violet,text width=3.4cm] (b1) at (8.8,2.15) {\textbf{$D$：106--150}\\窗口到期，deposit-only};
\node[box,fill=forceB,draw=rose,text width=3.4cm] (b2) at (14.3,2.15) {\textbf{$D$：151--260}\\RESUME(260) 立即补齐};
\node[box,kept,fill=settled,text width=2.6cm] (b3) at (19.6,2.15) {\textbf{新 session}\\从 $D[260]$ 出 261};
\draw[ref] (b0.north east) -- (5.5,5.0) -- (bold.west);
\draw[flow,draw=violet] (b0.south east) -- (5.5,2.15) -- (b1.west);
\draw[recover] (b1) -- (b2); \draw[recover] (b2) -- (b3);
\node[note,anchor=west,text=green] at (.75,.75) {RESUME 后：$G$ 与 $D$ 在 105 及以前一致；106--260 采用 $D$ 已确定的块。};
\node[note,anchor=west] at (.75,-.1) {绿色粗框：保留的派生前缀 / 新起点；灰色虚框：放弃的旧后缀；紫色：窗口回退；粉色：override 填充。};
\end{tikzpicture}
